\documentclass[10pt]{beamer}

\usetheme[progressbar=frametitle]{metropolis}
\usepackage{appendixnumberbeamer}
\usepackage{xcolor}
\usepackage{booktabs}
\usepackage{amsmath}
\usepackage{graphicx}
\usepackage{xspace}
\usepackage{verbatim}

\newcommand{\themename}{\textbf{\textsc{metropolis}}\xspace}

\setbeamercolor{background canvas}{bg=white}

\title{De-censoring Top-coded Wages in Austrian Social Security Data}
\subtitle{Initial Master Thesis Presentation}
\author{Felix Kurz \\Supervisors: 
Rudolf Winter-Ebmer $\And$ Wolfgang Frimmel}
\date{\today}

\begin{document}

\frame{\titlepage}

%--------------------------------------------------
\begin{frame}{Motivation}
Administrative data are central to empirical labor economics, which by its nature relies heavily on wage information.

\vspace{0.3cm}
Main advantages:
\begin{itemize}
    \item very large samples,
    \item reliable earnings records,
    \item long worker-firm histories.
\end{itemize}

\vspace{0.3cm}
Main limitation:
\begin{itemize}
    \item social security wage records are top-coded for high earners in Austria,
    \item observed contribution bases then differ from true wages.
\end{itemize}

\vspace{0.3cm}
This creates a problem for any analysis based on wages or on social security contribution bases as a proxy.
\end{frame}

%--------------------------------------------------
\begin{frame}{Why Does Censoring Matter?}
Ignoring top-coding can distort any analysis and insight drawn from the data.

\vspace{0.3cm}
To the extent that censored observations are included, results may not generalise to the full population.

\vspace{0.3cm}
The following illustrations show how censoring affects the observed wage distribution.
\end{frame}

%--------------------------------------------------
\begin{frame}{Illustration: Assumed Wage Distribution}
\centering
\includegraphics[width=0.85\textwidth]{iniPres/Plot_1.png}

\vspace{0.3cm}
Assume wages follow a standard normal distribution.
\end{frame}

%--------------------------------------------------
\begin{frame}{Illustration: The Top 10 Percent}
\centering
\includegraphics[width=0.85\textwidth]{iniPres/Plot_2.png}

\vspace{0.3cm}
The highest 10 percent of wages are affected by top-coding.
\end{frame}

%--------------------------------------------------
\begin{frame}{Illustration: What Censoring Does}
\centering
\includegraphics[width=0.85\textwidth]{iniPres/Plot_3.png}

\vspace{0.3cm}
All values above the ceiling are collapsed onto a single point.
\end{frame}

%--------------------------------------------------
\begin{frame}{Illustration: The Resulting Data}
\centering
\includegraphics[width=0.85\textwidth]{iniPres/Plot_4.png}

\vspace{0.3cm}
In practice, this creates a visible pile-up at the censoring threshold.
\end{frame}

%--------------------------------------------------
\begin{frame}{Research Question}
\begin{center}
Can top-coded wages in Austrian administrative data be predicted more accurately than with a classical Tobit-type correction model?
\end{center}
\end{frame}

%--------------------------------------------------
\begin{frame}{What This Thesis Does}
The thesis will:
\begin{enumerate}
    \item replicate and adapt an established Tobit imputation approach,
    \item use machine-learning methods as a flexible alternative,
    \item evaluate both against matched Lohnzettel wages as a benchmark.
\end{enumerate}

\vspace{0.3cm}
Additional question:
\begin{itemize}
    \item Can interpretable machine-learning methods provide additional insights beyond classical models?
\end{itemize}
\end{frame}

%--------------------------------------------------
% LITERATURE (moved before data, as requested)
%--------------------------------------------------
\begin{frame}{Literature Review}
This thesis builds on two strands of literature:
\begin{enumerate}
    \item \textbf{Card, Heining \& Kline (2013):} wage inequality using German administrative data
    \item \textbf{Gartner (2005):} imputation of censored wages in administrative data
\end{enumerate}

\vspace{0.3cm}
Their roles differ:
\begin{itemize}
    \item Card, Heining \& Kline provide the empirical framework and motivation.
    \item Gartner provides the direct methodological baseline.
\end{itemize}
\end{frame}

%--------------------------------------------------
% CARD SLIDE: rewritten to explain CHK (2013) specifically
%--------------------------------------------------
\begin{frame}{Literature I: Card, Heining \& Kline (2013)}
Card, Heining \& Kline study the rise of wage inequality in West Germany using linked employer--employee data from the IAB.

\vspace{0.3cm}
\begin{itemize}
    \item They decompose wages into worker effects, firm effects, and covariates (AKM model).
    \item The underlying administrative wage data are top-coded.
    \item To deal with censoring, they use a Tobit imputation procedure.
\end{itemize}

%\vspace{0.3cm}
%This is the starting point for this thesis: replicate and improve on their Tobit-based imputation using Austrian data and modern ML methods.
\end{frame}

%--------------------------------------------------
\begin{frame}{Literature II: Gartner (2005)}
Gartner provides the direct methodological foundation.

\vspace{0.3cm}
\begin{itemize}
    \item Administrative wage measures are institutionally determined.
    \item Contribution ceilings create top-coding and censoring.
    \item For high earners, observed wages are censored at the contribution ceiling and no longer reflect true earnings.
    \item For Germany, only the resulting Tobit-estimated wages are used in practice.
\end{itemize}

\vspace{0.3cm}
This creates a baseline to build on further and to test against.
\end{frame}

%--------------------------------------------------
% NEW SLIDE: What is new (professors requested this after literature)
%--------------------------------------------------
\begin{frame}{What Is New?}
This thesis extends the existing approach in two directions.

\vspace{0.3cm}
\textbf{1. Replication with Austrian data}
\begin{itemize}
    \item Adapt the Card, Heining \& Kline / Gartner Tobit imputation to Austrian administrative data.
\end{itemize}

\vspace{0.3cm}
\textbf{2. Machine-learning alternative}
\begin{itemize}
    \item Because Lohnzettel (LZ) data are available, a supervised ML approach becomes possible.
    \item LZ wages serve as the training target and evaluation benchmark.
    \item This is the main contribution.
\end{itemize}
\end{frame}

%--------------------------------------------------
\begin{frame}{Expected Contribution}
This thesis aims to contribute in three dimensions:
\begin{enumerate}
    \item \textbf{Measurement:} improved wage information in censored administrative data.
    \item \textbf{Method:} systematic comparison of a structured econometric baseline and a flexible ML approach.
    \item \textbf{Application:} a corrected wage measure usable in future research on inequality, mobility, and firm effects.
\end{enumerate}
\end{frame}

%--------------------------------------------------
% DATA SLIDES (moved after contribution, as requested)
%--------------------------------------------------
\begin{frame}{Administrative Data Landscape}
The starting point is a rich set of Austrian administrative data.

\vspace{0.3cm}
The available sources include:
\begin{itemize}
    \item \textbf{BG}: social security contribution basis and employment information,
    \item \textbf{Stammdaten}: worker demographics (gender, birth year, nationality),
    \item \textbf{Qual}: spell information (live births), and worker qualifications,
    \item \textbf{DG data}: firm-level characteristics (region, industry).
\end{itemize}

\vspace{0.3cm}
These data are rich in worker, job, and firm information, but the wage measure itself is censored.
\end{frame}

%--------------------------------------------------
\begin{frame}{Why Lohnzettel Data Matter}
The availability of \textbf{Lohnzettel (LZ)} data changes the empirical setting.

\vspace{0.3cm}
LZ data provide:
\begin{itemize}
    \item a wage measure that is not subject to top-coding,
    \item matched BG-LZ observations for the same worker-firm-year.
\end{itemize}

\vspace{0.3cm}
This creates new opportunities:
\begin{itemize}
    \item not only to correct censored wages,
    \item but also to assess how well different methods recover true wages,
    \item and to build a model that can impute wages even when no LZ data are available.
\end{itemize}
\end{frame}

%--------------------------------------------------
% EMPIRICAL APPROACH (renamed from "Implication", motivation text cut)
%--------------------------------------------------
\begin{frame}{Empirical Approach}
The thesis will proceed in three steps:
\begin{enumerate}
    \item Adapt a Tobit imputation model as a structured baseline.
    \item Train an XGBoost model as a flexible ML alternative.
    \item Evaluate both against matched LZ wages.
\end{enumerate}
\end{frame}

%--------------------------------------------------
% TOBIT METHODOLOGY (NEW, with equation from code)
%--------------------------------------------------
\begin{frame}{Methodology I: Tobit Imputation}
The baseline follows a Tobit-type interval regression.

\vspace{0.2cm}
\textbf{Model:}
\[
\ln w_{it}^{*} = \mathbf{x}_{it}' \boldsymbol{\beta} + \varepsilon_{it}, \quad \varepsilon_{it} \sim N(0, \sigma^2)
\]

For censored observations, only $\ln w_{it}^{*} \in [\text{observed}, \; \text{ceiling}]$ is known. %//vielleicht ein bisserl zu theoretisch?

\vspace{0.2cm}
\textbf{Covariates} $\mathbf{x}_{it}$:
\begin{columns}[T]
\begin{column}{0.48\textwidth}
\begin{itemize}
    \item age
    \item leave-one-out mean wage
    \item leave-one-out censoring share
    \item firm size, firm size$^2$
\end{itemize}
\end{column}
\begin{column}{0.48\textwidth}
\begin{itemize}
    \item coworker mean wage
    \item coworker censoring share
    \item one-year, one-person, small-firm dummies
\end{itemize}
\end{column}
\end{columns}

\vspace{0.2cm}
Estimated separately by \textbf{gender} $\times$ \textbf{age group}, year by year.
\end{frame}

%--------------------------------------------------
% XGBOOST METHODOLOGY (NEW, with inputs from code)
%--------------------------------------------------
\begin{frame}{Methodology II: XGBoost}
The ML alternative uses gradient-boosted trees (XGBoost).

\vspace{0.2cm}
\textbf{Target:} LZ daily wage (log, CPI-adjusted to 2000 euros).

\vspace{0.2cm}
\textbf{Features} (derived from BG data only):
\begin{columns}[T]
\begin{column}{0.48\textwidth}
\begin{itemize}
    \item all Tobit covariates
    \item gender, birth year
    \item number of observed years
    \item censoring flag
\end{itemize}
\end{column}
\begin{column}{0.48\textwidth}
\begin{itemize}
    \item interval bounds
    \item year
    \item federal state (Bundesland)
    \item nationality proxy (STSL)
\end{itemize}
\end{column}
\end{columns}

\vspace{0.2cm}
\textbf{Setup:}
\begin{itemize}
    \item 80/20 train--test split,
    \item 5-fold cross-validation with early stopping.
\end{itemize}

%\vspace{0.2cm}
%Key advantage: the model is trained on BG features and can predict wages for future data \textbf{without} LZ.
\end{frame}

\begin{comment}
 %--------------------------------------------------
\begin{frame}{Conclusion}
Top-coding creates a relevant measurement problem in administrative wage data.

\vspace{0.3cm}
This thesis asks whether:
\begin{itemize}
    \item censored wages can be recovered more credibly,
    \item rich worker--firm information improves prediction,
    \item flexible methods improve on standard correction models.
\end{itemize}

\vspace{0.3cm}
The key empirical opportunity is the availability of matched LZ wages as a benchmark.
\end{frame}

%--------------------------------------------------   
\end{comment}



\begin{frame}{References}
\footnotesize
Card, D., Heining, J., \& Kline, P. (2013). Workplace heterogeneity and the rise of West German wage inequality. \emph{Quarterly Journal of Economics}, 128(3), 967--1015.

\vspace{0.15cm}
Gartner, H. (2005). \emph{The imputation of wages above the contribution limit with the German IAB employment sample}. FDZ Methodenreport, 2, 2005.

\end{frame}

\end{document}